\subsection{Collision Notes} \label{sec:col-notes}
NOTE:
\emph{This is a place for making notes about the collision integrals.
  Eventually it will replace or be combined with the previous section.}

We consider three classes of collision:
%
\begin{enumerate}
\item Electron-neutral elastic collisions, e.g., $\ce{ Ar + e -> Ar + e}$;
\item Electron-neutral excitation reactions, e.g., $\ce{ Ar + e <=> Ar^{\ast} + e}$;
\item Electron-neutral ionization reactions, e.g., $\ce{ Ar^{\ast} + e <=> Ar^+ + e + e}$.
\end{enumerate}
%

\subsubsection{Electron-Neutral Elastic Scattering}
Let $()_e$ denote electron quantities---e.g., $f_e$ is the electron
distribution function, $m_e$ is the mass of an electron---and let
$()_N$ denote the analogous quantities for the neutral particle.
Then, the net rate of production of particles (per unit volume in
state space) is given by
%
\begin{equation}
Q_{EL}(\mbf{v}_e)
=
\int_{\mbb{R}^3} \int_{\mbb{R}^3}\int_{\mbb{R}^3}
\left(
f_e(\mbf{v}_e^{\prime}) f_N(\mbf{v}_N^{\prime})
-
f_e(\mbf{v}_e) f_N(\mbf{v}_N)
\right) W_{EL}(\mbf{v}_e, \mbf{v}_N; \mbf{v}_e^{\prime}, \mbf{v}_N^{\prime})
\,
d^3\mbf{v}_N \, d^3\mbf{v}_e^{\prime} \, d^3\mbf{v}_N^{\prime}
\label{eqn:elastic-collision-9d}
\end{equation}
%
where $\mbf{v}_e$ and $\mbf{v}_N$ denote the ``pre-collision''
velocities, $\mbf{v}_e^{\prime}$ and $\mbf{v}_N^{\prime}$ denote the
``post-collision'' velocities, and $W_{EL}$ denotes the elastic
collision transition probability, which encodes both constraints
(conservation of momentum and energy) \emph{and} information about the
relative likelihood of allowable transitions.

The nine dimensional integral in~\eqref{eqn:elastic-collision-9d} can
be reduced to a five dimensional integral by imposing conservation of
momentum and energy.  To do so in the simplest possible way, let
$()_0$ denote the quantities associated with the center of mass of the
two-particle system.  Specifically, the pre-collision position and
velocity of the center of mass are given by
%
\begin{equation*}
\mbf{x}_0 = \frac{m_N \mbf{x}_N + m_e \mbf{x}_e}{m_N + m_e},
\quad
\mbf{v}_0 = \frac{m_N \mbf{v}_N + m_e \mbf{v}_e}{m_N + m_e}
\end{equation*}
%
Conservation of momentum requires that
%
\begin{gather*}
m_N \mbf{v}_N + m_e \mbf{v}_e = m_N \mbf{v}_N^{\prime} + m_e \mbf{v}_e^{\prime}\\
(m_N +m_e) \mbf{v}_0 = (m_N + m_e) \mbf{v}_0^{\prime}.
\end{gather*}
%
That is, the velocity of the center of mass is not altered by the
collision.  Thus, letting $\mbf{u}$ denote velocities measured in a
frame attached to the center of mass, one may write
%
\begin{equation*}
m_N \mbf{u}_N^{\prime} + m_e \mbf{u}_e^{\prime}
=
m_N \mbf{v}_N^{\prime} + m_e \mbf{v}_e^{\prime} - (m_N + m_e) \mbf{v}_0
= 0.
\end{equation*}
%
Thus, conservation of momentum implies the following relationship
between the post-collision relative velocities:
%
\begin{equation*}
\mbf{u}_N^{\prime} = -\,\frac{m_e}{m_N} \mbf{u}_e^{\prime}.
\end{equation*}
%
To continue, let $\mbf{\hat{e}}$ a unit vector in the direction of the
relative velocity between the post-collision particles.  Then, the
relative post-collision velocity may be written as
%
\begin{equation*}
\mbf{u}_N^{\prime} - \mbf{u}_e^{\prime} = g^{\prime} \mbf{\hat{e}},
\end{equation*}
%
where $g'$ denotes the the magnitude of the post-collision relative
velocity.  To continue, we use conservation of energy to determine
$g^{\prime}$.

In the frame moving with the center of mass, conservation of energy requires that
%
\begin{equation*}
\frac{1}{2} m_N |\mbf{u}_N|^2 + \frac{1}{2} m_e |\mbf{u}_e|^2
=
\frac{1}{2} m_N |\mbf{u}_N^{\prime}|^2 + \frac{1}{2} m_e |\mbf{u}_e^{\prime}|^2.
\end{equation*}
%
In order to solve for $g^{\prime}$, we express the post-collision
velocities in terms of $g^{\prime}$ and $\mbf{\hat{e}}$.  By
definitions and previous results, one can show that
%
\begin{gather*}
\mbf{u}_N^{\prime} = \frac{m_e}{m_N + m_e} g^{\prime} \mbf{\hat{e}},\\
\mbf{u}_e^{\prime} = \frac{-m_N}{m_N + m_e} g^{\prime} \mbf{\hat{e}}.
\end{gather*}
%
Substituting into conservation of energy gives
%
\begin{equation*}
\frac{1}{2} m_N |\mbf{u}_N|^2 + \frac{1}{2} m_e |\mbf{u}_e|^2
=
\frac{1}{2} m_N \frac{m_e^2}{(m_N+m_e)^2} (g^{\prime})^2
+
\frac{1}{2} m_e \frac{m_N^2}{(m_N+m_e)^2} (g^{\prime})^2
=
\frac{1}{2} \frac{m_e m_N}{m_e + m_N} (g^{\prime})^2.
\end{equation*}
%
Thus,
%
\begin{equation*}
g^{\prime} = \left( \frac{m_e + m_N}{m_e m_N} \left( m_N |\mbf{u}_N|^2 + m_e |\mbf{u}_e|^2 \right) \right)^{1/2}.
\end{equation*}
%
This expression can be further simplified by introducing the
pre-collision relative speed $g$:
%
\begin{equation*}
g = | \mbf{u}_N - \mbf{u}_e |
\end{equation*}
%
Then,
\begin{equation*}
\frac{1}{2} m_N |\mbf{u}_N|^2 + \frac{1}{2} m_e |\mbf{u}_e|^2
=
\frac{1}{2} \frac{m_e m_N}{m_e + m_N} g^2,
\end{equation*}
%
which implies that $g^{\prime} = g$.

Thus, given the pre-collision velocities $\mbf{v}_N$ and $\mbf{v}_e$
and the post-collision relative velocity direction $\mbf{\hat{e}}$,
the post-collision velocities are completely determined:
%
\begin{gather*}
\mbf{v}_N^{\prime}
=
\mbf{v}_0 + \mbf{u}_N^{\prime}
=
m_N \mbf{v}_N + m_e \mbf{v}_e + \frac{m_e}{m_N+m_e} g \mbf{\hat{e}},\\
%
\mbf{v}_N^{\prime}
=
\mbf{v}_0 + \mbf{u}_N^{\prime}
=
m_N \mbf{v}_N + m_e \mbf{v}_e - \frac{m_N}{m_N+m_e} g \mbf{\hat{e}}.
\end{gather*}
%
Using these constraints, the collision integral
in~\eqref{eqn:elastic-collision-9d} may be rewritten as follows:
%
\begin{equation}
Q_{EL}(\mbf{v}_e)
=
\int_{\mbb{R}^3} \int_{S^2}
\left(
f_e(\mbf{v}_e^{\prime}) f_N(\mbf{v}_N^{\prime})
-
f_e(\mbf{v}_e) f_N(\mbf{v}_N)
\right) B(\mbf{v}_e, \mbf{v}_N; \mbf{\hat{e}})
\,
d\mbf{\hat{e}} \, d^3\mbf{v}_N,
\label{eqn:elastic-collision-5d}
\end{equation}
%
where we integrate only over transitions allowed by conservation of
momentum and energy and $B(\mbf{v}_e, \mbf{v}_N; \mbf{\hat{e}})$ has
taken the place of $W_{EL}$ for these allowable transitions.  It is
common to formulate $B$ as follows
%
\begin{equation*}
B(\mbf{v}_e, \mbf{v}_N; \mbf{\hat{e}}) = g \sigma(g,\mbf{\hat{e}}),
\end{equation*}
%
where $\sigma$ is the differential collision cross section.
Substituting this into~\eqref{eqn:elastic-collision-5d} gives the
typical form:
%
\begin{equation}
Q_{EL}(\mbf{v}_e)
=
\int_{\mbb{R}^3} \int_{S^2}
\left(
f_e(\mbf{v}_e^{\prime}) f_N(\mbf{v}_N^{\prime})
-
f_e(\mbf{v}_e) f_N(\mbf{v}_N)
\right) g \sigma(g, \mbf{\hat{e}})
\,
d\mbf{\hat{e}} \, d^3\mbf{v}_N,
\label{eqn:elastic-collision-5d}
\end{equation}
%
where, when evaluating the integral (e.g., using quadrature), the
quantities $\mbf{v}_e^{\prime}$, $\mbf{v}_N^{\prime}$, and $g$ are
determined from $\mbf{v}_e$, $\mbf{v}_N$, and $\mbf{\hat{e}}$ as
discussed previously.



\subsubsection{Electron-Neutral Excitation Reactions}


\subsubsection{Electron-Neutral Ionization Reactions}

